\begin{proof}[Proof of \cref{obj:lem:bv-maximal-cell-paths}]
\noindent\textit{1. Set up the count law and path envelopes.}
Let \(\mathbb P_{\mathrm{id}}\) be the joint law of the pilot and evaluation count tables in \cref{obj:def:process-handle}. For a count table \(\omega\), write \(W_j(\omega)\) for the path in the statement and set
\[
\rho_j(\omega)
=\|W_j(\omega)\|_\infty+
  \operatorname{TV}_{[0,1]}(W_j(\omega)).
\]
The second-moment assumption makes both \(\rho_j^2\) and
\(\|W_j\|_\infty^2\) integrable. The regularity assumptions give
Bochner-integrable paths, and the independence assumption applies to their
joint law. These observations also cover \(d=0\), when every sum below is
empty. % lean: CausalSmith.Stat.DiscreteBudgetvalueCurve.idealPilotErrorCellPath_centered_maximal

\noindent\textit{2. Reduce the centered sum to signed differences.}
Take independent count tables \(\omega,\omega'\) with law
\(\mathbb P_{\mathrm{id}}\), and define, for \(j=1,\ldots,d\) and
\(\lambda\in[0,1]\),
\[
\Delta_j(\lambda)
=W_j(\lambda;\omega)-W_j(\lambda;\omega').
\]
Let \(\eta_1,\ldots,\eta_d\) be independent uniform signs, independent of
the count tables, and define their signed supremum energy by
\[
\mathcal E_\eta(\Delta)
=\mathbb E_\eta\left\|
  \sum_{j=1}^d\eta_j\Delta_j
 \right\|_\infty^2.
\]
Jensen's inequality applied to the independent copy gives
\[
\mathbb E\left\|
 \sum_{j=1}^d(W_j-\mathbb EW_j)
 \right\|_\infty^2
\le
\mathbb E_{\omega,\omega'}\left\|
 \sum_{j=1}^d\Delta_j
 \right\|_\infty^2.
\]
The pairs \((W_j(\omega),W_j(\omega'))\) are independent across \(j\).
Swapping either entry of any pair preserves their joint law and reverses
the sign of its difference. Averaging these equalities over all sign
choices therefore yields
\[
\mathbb E_{\omega,\omega'}\left\|
 \sum_{j=1}^d\Delta_j
 \right\|_\infty^2
=
\mathbb E_{\omega,\omega'}\mathcal E_\eta(\Delta).
\]
All these expectations are finite by the second-moment assumption.
% lean: Causalean.Stat.Concentration.BoundedVariation.independent_copy_difference_energy_eq_sign_energy

\noindent\textit{3. Bound the signed supremum for fixed count tables.}
Fix a pair \((\omega,\omega')\) for which every \(W_j\) has finite
variation. Let \(\mathcal V_j(t)\), \(t\in[0,1]\), be the cumulative
variation of \(\Delta_j\) from \(0\) to \(t\). Continuity and finite
variation give a continuous, nondecreasing control satisfying
\[
\mathcal V_j(0)=0,\qquad
\mathcal V_j(1)=\operatorname{TV}_{[0,1]}(\Delta_j),\qquad
|\Delta_j(t)-\Delta_j(s)|
 \le \mathcal V_j(t)-\mathcal V_j(s)\quad(s\le t).
\]
Define the common control on \([0,1]\) by
\[
\mathcal U(t)=\sum_{j=1}^d
  \mathcal V_j(1)\mathcal V_j(t).
\]
Since \(0\le\mathcal V_j(t)-\mathcal V_j(s)\le\mathcal V_j(1)\),
\begin{equation}
\sum_{j=1}^d\bigl(\Delta_j(t)-\Delta_j(s)\bigr)^2
 \le \mathcal U(t)-\mathcal U(s)
 \quad(s\le t).
\label{eq:obj:bv-maximal-cell-paths-control}
\end{equation}

Choose ordered, compatible dyadic control quantiles \(t_{k,i}\), for
\(k\ge0\) and \(0\le i\le2^k\), so that
\[
\mathcal U(t_{k,i})=\frac{i}{2^k}\mathcal U(1).
\]
Such choices follow from continuity of \(\mathcal U\). For each sign
vector, let
\[
M_k(\eta)=
\max_{0\le i<2^{k+1}}
\left|
\sum_{j=1}^d\eta_j
 \bigl(\Delta_j(t_{k+1,i+1})-\Delta_j(t_{k+1,i})\bigr)
\right|.
\]
By \cref{eq:obj:bv-maximal-cell-paths-control}, every edge at level \(k+1\) has squared coefficient sum at most
\(\mathcal U(1)/2^{k+1}\). For \(k\ge0\), \(0\le i<2^{k+1}\), and \(1\le j\le d\), write
\[
\gamma_{k,i,j}
=\Delta_j(t_{k+1,i+1})-\Delta_j(t_{k+1,i}),
\qquad
\mathcal Y_k(\eta)=M_k(\eta)^2,
\qquad
\mathfrak v_k=\frac{\mathcal U(1)}{2^{k+1}}.
\]
When \(\mathfrak v_k>0\), each edge satisfies
\(\sum_j\gamma_{k,i,j}^2\le\mathfrak v_k\). For every
\(\chi>0\), independence of the signs and the elementary bound
\(\cosh z\le e^{z^2/2}\) give
\[
\mathbb E_\eta\exp\!\left(\chi\sum_j\eta_j\gamma_{k,i,j}\right)
=\prod_j\cosh(\chi\gamma_{k,i,j})
\le\exp\!\left(\frac{\chi^2\mathfrak v_k}{2}\right).
\]
Applying this bound to both signs, choosing
\(\chi=\sqrt{\mathfrak s}/\mathfrak v_k\), and taking the union over
\(2^{k+1}\) edges gives, for \(\mathfrak s>0\),
\[
\mathbb P_\eta\{\mathcal Y_k\ge\mathfrak s\}
\le 2^{k+2}\exp\!\left(-\frac{\mathfrak s}{2\mathfrak v_k}\right).
\]
Set \(\mathfrak t_k=2\mathfrak v_k\log(2^{k+2})\). Integrating the
tail above \(\mathfrak t_k\) yields
\[
\mathbb E_\eta\mathcal Y_k
\le\mathfrak t_k+
 \int_{\mathfrak t_k}^\infty
 2^{k+2}e^{-\mathfrak s/(2\mathfrak v_k)}\,d\mathfrak s
=2\mathfrak v_k\{\log(2^{k+2})+1\}
\le32\bigl(1+\log(2^{k+1})\bigr)\mathfrak v_k.
\]
When \(\mathfrak v_k=0\), \cref{eq:obj:bv-maximal-cell-paths-control}
makes every edge increment zero. Thus in every case
\begin{equation}
\mathbb E_\eta M_k(\eta)^2
\le
32\bigl(1+\log(2^{k+1})\bigr)
   \frac{\mathcal U(1)}{2^{k+1}}.
\label{eq:obj:bv-maximal-cell-paths-level}
\end{equation}
The numerical factors in \cref{eq:obj:bv-maximal-cell-paths-level}
have a summable square root:
\begin{equation}
\sum_{k=0}^\infty
 \sqrt{\frac{32(1+\log(2^{k+1}))}{2^{k+1}}}
\le
\sum_{k=0}^\infty 8\left(\frac34\right)^k
=32\le40.
\label{eq:obj:bv-maximal-cell-paths-series}
\end{equation}
Here the termwise bound follows from \(\log 2\le1\) and
\(k+2\le4(9/8)^k\).

To reconstruct the signed path, for each \(t\in[0,1]\) and
\(k\ge0\) define
\[
\xi(t)=
\begin{cases}
\mathcal U(t)/\mathcal U(1),&\mathcal U(1)>0,\\
0,&\mathcal U(1)=0,
\end{cases}
\qquad
\iota_k(t)=\min\{2^k-1,\lfloor2^k\xi(t)\rfloor\}.
\]
These lower dyadic indices start at zero, satisfy
\(\iota_{k+1}(t)\in\{2\iota_k(t),2\iota_k(t)+1\}\), and obey
\[
\frac{\iota_k(t)}{2^k}\le\xi(t)
\le\frac{\iota_k(t)+1}{2^k}.
\]
The parent grid point equals the even child by compatibility; the odd
child differs by one adjacent edge. At level zero the grid point and
\(0\) have the same control value, so \cref{eq:obj:bv-maximal-cell-paths-control}
makes their path coordinates equal. Telescoping therefore gives, for
every \(k\ge0\),
\[
\left|\sum_{j=1}^d\eta_j
 \Delta_j(t_{k,\iota_k(t)})\right|
\le
\left|\sum_{j=1}^d\eta_j\Delta_j(0)\right|
+\sum_{\nu=0}^{k-1}M_\nu(\eta),
\]
where the sum is empty at \(k=0\). The control values of
\(t_{k,\iota_k(t)}\) and \(t\) differ by at most
\(\mathcal U(1)/2^k\). Applying
\cref{eq:obj:bv-maximal-cell-paths-control} in whichever order these
two times occur gives
\[
\sum_{j=1}^d
 \bigl(\Delta_j(t_{k,\iota_k(t)})-\Delta_j(t)\bigr)^2
\le \frac{\mathcal U(1)}{2^k}.
\]
Hence the finite signed sums along the chain converge to their value
at \(t\), including on flat parts of the control. If
\(\mathcal U(1)=0\), the same control bound makes all \(\Delta_j\)
constant and the conclusion is immediate. Passing to the limit and
taking the supremum over \(t\) gives
\[
\left\|\sum_{j=1}^d\eta_j\Delta_j\right\|_\infty
\le
\left|\sum_{j=1}^d\eta_j\Delta_j(0)\right|
 +\sum_{k=0}^\infty M_k(\eta).
\]
Sign orthogonality, the \(L_2\) triangle inequality, and
\cref{eq:obj:bv-maximal-cell-paths-level,eq:obj:bv-maximal-cell-paths-series}
imply
\[
\mathcal E_\eta(\Delta)
\le
4096\left\{
 \sum_{j=1}^d\Delta_j(0)^2+\mathcal U(1)
\right\}
\le
4096\sum_{j=1}^d
 \left(\|\Delta_j\|_\infty+
 \operatorname{TV}_{[0,1]}(\Delta_j)\right)^2.
\]
Finally, the triangle inequalities for the supremum norm and total
variation give
\[
\mathcal E_\eta(\Delta)
\le
8192\sum_{j=1}^d
 \bigl(\rho_j(\omega)^2+\rho_j(\omega')^2\bigr).
\]
% lean: Causalean.Stat.Concentration.BoundedVariation.rademacherEnergy_le_pathSize_sq

\noindent\textit{4. Average the pointwise bound.}
Integrating the last inequality over the two independent count tables
and using their identical laws gives
\[
\mathbb E_{\omega,\omega'}\mathcal E_\eta(\Delta)
\le
8192\sum_{j=1}^d
 \mathbb E_{\omega,\omega'}
 \bigl(\rho_j(\omega)^2+\rho_j(\omega')^2\bigr)
=
16384\sum_{j=1}^d\mathbb E\rho_j^2.
\]
Combining this with the two relations in Step 2 proves the stated
inequality. % lean: Causalean.Stat.Concentration.BoundedVariation.centered_path_sum_maximal
\end{proof}
